\section{Approach}
\label{sec:approach}
To address these objectives, we design a compact Multilayer Perceptron (MLP) architecture (Dense $9 \to 8 \to 1$ with exactly 89 parameters) tailored as a \gls{tinyml} model. This shallow configuration is selected to balance representation capacity against the strict compute, memory, and transmission limits of edge devices, as mathematically and physically justified in Chapter 4. A detailed structural schematic of this 9:8:1 neural network topology, detailing layer configurations, normalization, and bias terms, is referenced in Chapter~\ref{ch:methodology} (corresponding to the asset \texttt{fig\_nn\_architecture1.png}). The 9 input features fed into this model consist of:
\begin{enumerate}
  \item \textbf{Spatial characteristics:} Log-distance between the transmitter and receiver, the count of concrete/brick walls (\(W_{\text{brick}}\)), and the count of wooden walls (\(W_{\text{wood}}\)).
  \item \textbf{Environmental covariates:} CO\textsubscript{2} levels, Relative Humidity, Particulate Matter (PM2.5), Barometric Pressure, and Temperature.
  \item \textbf{Link quality indicator:} Receiver-reported Signal-to-Noise Ratio (\gls{snr}).
\end{enumerate}

We construct a software simulation using python and tensorflow to run federated learning experiments over the 2,079,534 real-world measurements collected during the prior Hölderlinstraße campus campaign. To model client drift realistically, we partition the dataset by physical device location, establishing six virtual clients representing distinct indoor environments.

To combat the severe client drift arising from this non-\gls{iid} partitioning, we implement and evaluate the standard Federated Averaging (\gls{fedavg}) algorithm as a baseline against three advanced federated optimization algorithms:
\begin{itemize}
  \item \textbf{FedProx}, which adds a proximal regularization term to the local loss function to prevent local updates from straying too far from the global parameters.
  \item \textbf{SCAFFOLD}, which utilizes local and global control variates to correct the client-side gradient direction during local training.
  \item \textbf{FedYogi}, which applies a server-side adaptive optimizer to smooth the updates using momentum and variance estimation.
\end{itemize}
These federated models are compared against two key benchmarks: an idealized centralized MLP model (serving as the theoretical accuracy upper bound) and a structure-only multiple linear regression baseline (representing the traditional distance-and-wall propagation modeling standard).

We evaluate the models using three primary dimensions: path loss prediction accuracy (measured via $R^2$ and root mean square error (RMSE) in dB), convergence rate (number of communication rounds required to stabilize), and network communication overhead (measured by cumulative upload bytes per client). 

Finally, our approach relies on a hybrid methodology: we use rigorous python simulations to validate convergence and accuracy, and then combine this with memory and communication payload calculations to prove the physical feasibility of future on-device deployment. Specifically, we evaluate the Flash memory footprint, static RAM (\gls{sram}) requirements, and transmission Time-on-Air (\gls{toa}) payloads of our 89-parameter model against the hardware limitations of the Arduino MKR~WAN~1310.
